\section*{Περίληψη}


\noindent
Η συμπεριφορά των ελαστικών αποτελεί έναν από τους πιο καθοριστικούς παράγοντες επίδοσης και στρατηγικής στις σύγχρονες
αγωνιστικές διοργανώσεις, ενώ επίσης η πρόσφυσή τους εξαρτάται έντονα από τη θερμοκρασία και τη φθορά του
πέλματος. Ωστόσο, οι προσομοιωτές χρόνου γύρου χρησιμοποιούν
συνήθως σταθερές και ονομαστικές συνθήκες θερμοκρασίας και
φθοράς, ενώ τα μοντέλα θερμοκρασίας και φθοράς της βιβλιογραφίας βασίζονται κυρίως σε μη-μόνιμους
προσομοιωτές, που συνοδεύονται από υψηλό υπολογιστικό κόστος και πολυπλοκότητα, και δεν είναι κατάλληλοι για όλα τα είδη προσομοιώσεων.\\[8pt]
\noindent
Στην παρούσα διπλωματική εργασία αναπτύσσεται ένα υπολογιστικό πλαίσιο πρόβλεψης της θερμοκρασίας και της φθοράς των ελαστικών ενός αγωνιστικού οχήματος.
Διατυπώνεται ένα
φυσικό θερμοδυναμικό μοντέλο του πέλματος και ένα μοντέλο φθοράς τριών μηχανισμών, τα οποία συνδέονται με
την πρόσφυση και με ένα τετράτροχο (\textit{two track}) μοντέλο οχήματος. Αναπτύσσεται επίσης ένα ψηφιακό δίδυμο του οχήματος, μέσω εκτίμησης της κατάστασής του, δηλαδή των μεγεθών ολίσθησης (\textit{slip angle, slip ratio})
και των φορτίων των τροχών, με δεδομένα μόνο το προφίλ ταχύτητας και τη γεωμετρία της πίστας.
Τα αποτελέσματα του ψηφιακού διδύμου επιτρέπουν την προσέγγιση της θερμοκρασίας και της φθοράς κάθε τροχού από πραγματικά δεδομένα οχήματος. Με τη βοήθεια αυτού του μοντέλου, και σε συνδυασμό με τις αντίστοιχες μετρήσεις θερμοκρασίας και φθοράς, μπορεί να βαθμονομηθεί και έπειτα να επαληθευτεί το θερμοδυναμικό μοντέλο.\\[8pt]
Επιπλέον, κατασκευάζεται ένα χρονο-μεταβαλλόμενο διάγραμμα επιδόσεων (\textit{GGV diagram}), το οποίο ενσωματώνεται σε
υφιστάμενο προσομοιωτή ανοιχτού κώδικα ψευδο-μόνιμης κατάστασης (\textit{Quasi-Steady-State -- QSS}),
σε αντίθεση με θερμοδυναμικά μοντέλα της βιβλιογραφίας που υλοποιούνται σε μη-μόνιμους (\textit{transient}) προσομοιωτές. Οι μη-μόνιμοι προσομοιωτές, αν και ακριβέστεροι, είναι λιγότερο κατάλληλοι για μελέτες σχεδιασμού του οχήματος με μεγάλο αριθμό μεταβλητών σχεδιασμού, καθώς το υψηλό υπολογιστικό κόστος περιορίζει το πλήθος των προσομοιώσεων, ενώ το μοντέλο οδηγού εισάγει μεταβλητότητα στα αποτελέσματα που δεν οφείλεται στο ίδιο το όχημα.
\\[8pt]
Η μεθοδολογία εφαρμόζεται σε όχημα τύπου \textit{Formula 1} στην πίστα \textit{Circuit de
Barcelona-Catalunya}. Σε σύγκριση με πειραματικά δεδομένα της βιβλιογραφίας, το μοντέλο
αποτυπώνει ποιοτικά σωστά την εξέλιξη της θερμοκρασίας. Η επίδραση της θερμοκρασίας στον χρόνο γύρου δεν
είναι αμελητέα, καθώς ένας γύρος με ψυχρά ελαστικά προκύπτει περίπου $3\,\mathrm{s}$ βραδύτερος. Επιπλέον,
ο ρυθμός φθοράς διαφοροποιείται έντονα κατά μήκος της πίστας και μεταξύ των τροχών, ενώ μια ευρετική
μελέτη διατήρησης ελαστικών δείχνει πως μια πιο συντηρητική οδήγηση δεν μειώνει απαραίτητα τη φθορά,
καθώς παράλληλα μειώνεται η θερμοκρασία και κυριαρχεί η φθορά λόγω \textit{graining}, κάνοντας εμφανή την ανάγκη χρήσης αλγορίθμου βελτιστοποίησης με σκοπό τη στοχευμένη διατήρηση ελαστικών. Συμπεραίνεται πως η κατάσταση των ελαστικών μπορεί να ενσωματωθεί με απλό
τρόπο σε προσομοιωτή ψευδο-μόνιμης κατάστασης, διατηρώντας το χαμηλό υπολογιστικό κόστος και την
ανεξαρτησία από μοντέλο οδηγού.\\[8pt]
\noindent
Τα αποτελέσματα ερμηνεύονται κυρίως ποιοτικά, καθώς η βαθμονόμηση του μοντέλου περιορίστηκε από τα διαθέσιμα πειραματικά δεδομένα, ενώ η απλοποιημένη θερμική μοντελοποίηση του ελαστικού δεν αποτυπώνει πλήρως τη χρονική απόκρισή του. Ως μελλοντική κατεύθυνση, είναι απαραίτητη η βαθμονόμηση και επαλήθευση με μεγαλύτερο σύνολο από πειραματικά δεδομένα, ενώ προτείνεται και ο εμπλουτισμός του θερμικού μοντέλου με περισσότερα φαινόμενα, όπως η επίδραση ολόκληρης της θερμικής μάζας του τροχού και η παραγωγή θερμότητας από το σύστημα πέδησης. Επιπλέον, με τη συμπερίληψη των γωνιών ευθυγράμμισης των τροχών (\textit{toe, camber}), το πλαίσιο θα μπορεί να αξιοποιηθεί ποσοτικά στον σχεδιασμό και τη ρύθμιση (\textit{setup}) του οχήματος αλλά και της στρατηγικής αγώνα.